% Judul penuh meluber 26,51 pt di kepala halaman; hanya mark header yang
% dipendekkan (sekali pakai, lihat \judulpendek di preamble-en.tex).
\judulpendek{Estimation, Confidence Intervals, and Testing}
\chapter{Estimation, Confidence Intervals, and Hypothesis Testing}

A team tests an attack detector on $50$ labeled examples, and the detector catches $33$ of
them. The report written that night reads: recall $66$ percent. A second detector, tested on
$100$ examples, catches $53$: recall $53$ percent. The morning meeting decides to use the
first, because $66$ is larger than $53$, and that decision is finished in one minute.

The arithmetic is correct: $33/50=0.66$. What is wrong is the assumption that this number
points at one number. If those fifty examples are replaced by fifty other examples from the
same collection, the result is no longer $33$; it wobbles, and the wobble is not small. After
two hours of writing the program, the team is in fact holding two sentences of very different
strength. The first sentence, ``on this data we see $66$ percent'', says almost nothing. The
second, ``the true catch rate lies somewhere inside this interval'', says something usable,
provided we know how wide that interval is and what that width means.

Chapter 19 already gave two tools: the law of large numbers, which says the sample mean
concentrates at the true value, and the central limit theorem, which says the error spreads
normally with standard deviation $\sigma/\sqrt n$. Both are statements about many trials at
once, and neither answers the question of the morning meeting: from one trial already
finished, what sentence may we utter about a parameter we have never seen, and how often is
that sentence wrong if the procedure is repeated. This chapter builds the answer, from
estimators to p-values.

\begin{untukanak}
A farmer wants to know how many weeds there are in his field. The whole field is too large to
inspect, so he takes twenty small plots spread around, then counts the weeds in each plot.
From those twenty numbers he guesses the state of the whole field. The trouble is that a guess
is always a single number, whereas the state of a field is never that simple.

Now the part that makes one curious. In those twenty plots he finds no weeds at all. May he go
home and say that his field is clean? Certainly not. It may be that there are really no weeds,
and it may be that he merely happened to choose plots made of bare soil. All he can say is:
if the weeds were plentiful, twenty empty plots would be almost impossible. From that he can
compute an honest upper bound, for instance ``at most three in a hundred plots''.

The same holds when a teacher grades a few students' quizzes and when a computer program is
tested on a few examples.
In this chapter we shall learn why one number is never enough, why the honest answer is always
an interval, and why the interval widens when the evidence is thin.
\end{untukanak}

\section{Estimators, bias, and standard error}
\tingkat{basic}

Every statistical conclusion begins with one small decision: which quantity we compute from
the data. The mean, a proportion, the sample standard deviation, all are functions of the data
and require no knowledge of the parameter being sought. In everyday language: an estimator is
a surveyor who guesses the height of a tree from the length of its shadow, and we need to know
whether that surveyor is systematically off or merely unsteady.

\kamus{estimator}{penaksir}{a surveyor who guesses the height of a tree from the length of
its shadow, rather than from the tree itself}

\begin{definisi}[Statistic and estimator]
\index{estimator}
Let $X_1,\dots,X_n$ be a random sample from a distribution that depends on a parameter
$\theta$. A \emph{statistic} is a function of $X_1,\dots,X_n$ that does not contain
$\theta$. A statistic used to guess $\theta$ is called an \emph{estimator} and written
$\hat\theta$; the value obtained on one sample is called an \emph{estimate}.
\end{definisi}

An estimator is a random variable, because it is a function of random variables. Hence it has
an expectation and a variance, and those two quantities determine its quality. The expectation
answers the question ``does the guess keep missing in one direction'', while the variance
answers ``how far does the guess wobble from sample to sample''.

\begin{definisi}[Bias and unbiased estimator]
\index{unbiased}
The bias of the estimator $\hat\theta$ for $\theta$ is $\operatorname{bias}(\hat\theta)
=\mathbb{E}[\hat\theta]-\theta$. An estimator is called \emph{unbiased} if its bias is zero,
namely $\mathbb{E}[\hat\theta]=\theta$ for every possible value of $\theta$.
\end{definisi}

The sample mean is unbiased for $\mu$, and that was already used in the proof of the law of
large numbers in Chapter 19, so we do not repeat it here. More interesting is the estimator
used for the variance. With divisor $n$, the estimator
$S_n^2=\frac1n\sum(X_i-\bar X_n)^2$ is biased downward; with divisor $n-1$, the estimator
$S^2=\frac1{n-1}\sum(X_i-\bar X_n)^2$ is unbiased. The reason is that $\sum(X_i-\bar X_n)^2$
loses one degree of freedom because $\bar X_n$ is itself computed from the same data, so its
expectation is $(n-1)\sigma^2$ and not $n\sigma^2$. The divisor $n-1$ is a correction for that
loss, not mere tidiness.

The quality of an estimator is usually summarized by one number, the mean squared error. In
everyday language: how large an error we expect, computed in a way that punishes large errors
more heavily.

\begin{definisi}[Mean squared error]
\index{mean squared error}
$\operatorname{MSE}(\hat\theta)=\mathbb{E}[(\hat\theta-\theta)^2]$.
\end{definisi}

\begin{teorema}[MSE decomposition]
For an estimator $\hat\theta$ with finite variance,
\[
\operatorname{MSE}(\hat\theta)=\operatorname{bias}(\hat\theta)^2
+\operatorname{Var}(\hat\theta).
\]
\end{teorema}

\begin{proof}[Proof]
Write $\hat\theta-\theta=(\hat\theta-\mathbb{E}\hat\theta)
+(\mathbb{E}\hat\theta-\theta)$. Square it and take the expectation:
\[
\mathbb{E}[(\hat\theta-\theta)^2]
=\mathbb{E}[(\hat\theta-\mathbb{E}\hat\theta)^2]
+2(\mathbb{E}\hat\theta-\theta)\,\mathbb{E}[\hat\theta-\mathbb{E}\hat\theta]
+(\mathbb{E}\hat\theta-\theta)^2 .
\]
The middle term vanishes because $\mathbb{E}[\hat\theta-\mathbb{E}\hat\theta]=0$, while the
first term is the variance and the third is the bias squared.
\end{proof}

That identity explains why ``unbiased'' is not a championship title. An unbiased estimator may
have a large variance, and a slightly biased estimator may have a far smaller variance so that
its MSE is lower. To estimate $\mu$ from $n$ observations, the statistic $X_1$ is unbiased with
variance $\sigma^2$, whereas $\bar X_n$ is unbiased with variance $\sigma^2/n$; at $n=20$ the
difference is a factor of twenty.

The second example shows a bias that must be computed, because it is used every day for
proportions. If $K\sim\mathrm{Bin}(n,p)$ and $\hat p=K/n$, then
$\hat p(1-\hat p)$ is the natural estimator for $p(1-p)$, that is the variance of one trial.
Its expectation is
\[
\mathbb{E}[\hat p(1-\hat p)]
=\mathbb{E}[\hat p]-\mathbb{E}[\hat p^2]
=p-\Bigl(p^2+\frac{p(1-p)}{n}\Bigr)
=p(1-p)\Bigl(1-\frac1n\Bigr),
\]
so its bias is $-p(1-p)/n$. For $p=0.3$ and $n=20$ the value is $p(1-p)=0.21$
and the bias is $-0.0105$, that is five percent downward. At $n=200$ the bias is only
$-0.00105$. A bias that shrinks like $1/n$ is typical: it disappears at infinity, but it is
present in every finite sample.

The second property demanded of an estimator is that it actually heads for the target as the
data grow. In everyday language: with enough evidence, a wrong guess becomes increasingly
impossible, not merely increasingly rare.

\begin{definisi}[Consistent]
\index{consistent}
A sequence of estimators $\hat\theta_n$ is \emph{consistent} for $\theta$ if for every
$\varepsilon>0$ we have $P(|\hat\theta_n-\theta|\ge\varepsilon)\to0$ as
$n\to\infty$.
\end{definisi}

\begin{proposisi}[MSE going to zero suffices for consistency]
\index{consistent}
If $\operatorname{MSE}(\hat\theta_n)\to0$, then $\hat\theta_n$ is consistent for
$\theta$.
\end{proposisi}

\begin{proof}[Proof]
By Chebyshev's inequality, already proved in Chapter 19,
\[
P(|\hat\theta_n-\theta|\ge\varepsilon)
\le\frac{\mathbb{E}[(\hat\theta_n-\theta)^2]}{\varepsilon^2}
=\frac{\operatorname{MSE}(\hat\theta_n)}{\varepsilon^2}.
\]
For fixed $\varepsilon$ the right-hand side goes to zero, and the inequality closes the
proof.
\end{proof}

Note that Chebyshev uses the MSE, not the variance. The bias term is squared, so a biased
estimator is still consistent provided its bias goes to zero. Conversely, the estimator $X_1$
is unbiased for every $n$ but never uses the rest of the data, so it is not consistent.

The next quantity is used continually until the end of the chapter. The variance of an
estimator says how far the guess wobbles, and its square root has a name of its own.

\kamus{standard error}{galat baku}{the range of wobble of one guess, not the width of the whole
data set}

\begin{definisi}[Standard error]
\index{standard error}
The \emph{standard error} of the estimator $\hat\theta$ is $\operatorname{se}(\hat\theta)
=\sqrt{\operatorname{Var}(\hat\theta)}$. If the formula contains an unknown parameter, the
\emph{estimated} standard error is obtained by replacing that parameter with its estimate.
\end{definisi}

For a proportion, $\operatorname{se}(\hat p)=\sqrt{p(1-p)/n}$, and the estimated version is
$\sqrt{\hat p(1-\hat p)/n}$. For the data $33$ out of $50$, $\hat p=0.66$ and
\[
\sqrt{\frac{0.66\cdot0.34}{50}}=\sqrt{0.004488}=0.066993 .
\]
For twelve measurements of delay time (in milliseconds),
\begin{align*}
&18.2,\ 19.1,\ 17.8,\ 20.3,\ 18.9,\ 19.7,\\
&18.4,\ 21.0,\ 17.5,\ 19.9,\ 18.6,\ 20.1,
\end{align*}
the mean is $19.125$ and the sample standard deviation $1.082191$. The standard error is
$1.082191/\sqrt{12}=0.312401$.
That number does not say that the true mean is $0.312$ away from
$19.125$; it says that guesses like this one, if repeated, wobble on a scale of about
$0.31$. The confidence interval is the way to turn the last sentence into a sentence we can
use.

\section{Confidence intervals for a mean}
\tingkat{intermediate}

A single number without a measure of wobble is easily misused. What is needed is a pair of
numbers that, with high probability, traps the parameter. In everyday language: we do not
shoot at one point, we set a net and promise that the net catches the fish we are after on most
throws.

\kamus{confidence interval}{interval kepercayaan}{a net set with a certain width, and a promise
that the net catches the target on most throws}

\begin{definisi}[Confidence interval and coverage]
\index{confidence interval}
Let $C(X_1,\dots,X_n)$ be an interval computed from the data. $C$ is called a
\emph{confidence interval} of level $1-\alpha$ for $\theta$ if
$P_\theta(\theta\in C)\ge1-\alpha$ for every $\theta$. That probability is called the
\emph{coverage} of the procedure.
\end{definisi}

Note where the probability sits: $P_\theta$ is computed before the data are taken, and what is
random is $C$, not $\theta$. This is where misinterpretation most often occurs, and we return
there once the formula exists.

\begin{teorema}[The z-interval]
\index{confidence interval}
Let $X_1,\dots,X_n$ be independent and distributed $N(\mu,\sigma^2)$ with $\sigma$
known. For $z=z_{1-\alpha/2}$, the interval
\[
C=\Bigl[\bar X_n-z\frac{\sigma}{\sqrt n},\ \bar X_n+z\frac{\sigma}{\sqrt n}\Bigr]
\]
has coverage exactly $1-\alpha$. For other distributions with finite variance, its
coverage tends to $1-\alpha$ as $n\to\infty$.
\end{teorema}

\begin{proof}[Proof]
Let $Z=\sqrt n(\bar X_n-\mu)/\sigma$. For a normal sample, $Z\sim N(0,1)$ exactly;
for other distributions, $Z$ tends to $N(0,1)$ by the central limit theorem in Chapter 19.
The inequality $\mu\in C$ is equivalent to $|Z|\le z$, because
\[
|\bar X_n-\mu|\le z\frac{\sigma}{\sqrt n}
\iff \frac{\sqrt n\,|\bar X_n-\mu|}{\sigma}\le z .
\]
So $P(\mu\in C)=P(|Z|\le z)=1-\alpha$ in the normal case, and tends to $1-\alpha$
in the general case.
\end{proof}

If $\sigma$ is unknown, as in almost all real data, replacing
$\sigma$ by $s$ adds uncertainty and thickens the tails. The result uses the
$t$ distribution with $n-1$ degrees of freedom, published by Gosset in 1908 under the pen
name Student:
\[
C=\Bigl[\bar X_n-t_{n-1,1-\alpha/2}\frac{s}{\sqrt n},\
\bar X_n+t_{n-1,1-\alpha/2}\frac{s}{\sqrt n}\Bigr],
\]
with coverage exactly $1-\alpha$ if the sample is normal. The proof requires the joint
distribution of $\bar X_n$ and $s$ to be independent, and it is in Wasserman, \emph{All of
Statistics}, Chapter 6. For the $n=12$ delay-time data above, $t_{11,0.975}=2.200985$, so
the half-width is $2.200985\cdot0.312401=0.687591$ and the interval is
$(18.4374;\ 19.8126)$. With $z=1.96$ the half-width is $0.612296$: admitting that
$\sigma$ is unknown widens the interval by twelve percent at $n=12$, and by one
percent at $n=100$.

Now the interpretation, and here we must go slowly. Once the data are taken,
$C$ is no longer a random variable: it is the fixed interval $(18.4374;\ 19.8126)$. The
true mean $\mu$ is also a fixed number, only unknown. The sentence
``the probability that $\mu$ lies inside $(18.4374;\ 19.8126)$ equals $0.95$''
is therefore meaningless in this framework: no probability is left, because
nothing is random. What has probability $0.95$ is \emph{the procedure}, and that is a
statement about a trial not yet run.

Figure~\ref{fig:dua-puluh-selang} runs that trial twenty times, with
$\mu=19.125$ and $\sigma=1$ known: nineteen intervals contain $\mu$, one
does not.

\begin{figure}[htbp]
\centering
\begin{tikzpicture}[x=4.15cm,y=0.235cm,font=\scriptsize]
  \draw[->,layergrey] (17.80,0)--(20.22,0);
  \draw[layergrey,densely dashed,thick] (19.125,0)--(19.125,21.4);
  \node[above,font=\scriptsize] at (19.125,21.4) {$\mu=19.125$};
  \foreach \y/\a/\b in {
    1/18.5955/19.7271,
    2/18.5211/19.6527,
    3/18.7441/19.8757,
    4/18.5895/19.7211,
    5/18.4046/19.5362,
    6/18.6636/19.7952,
    7/18.9356/20.0672,
    8/18.8326/19.9642,
    9/18.3561/19.4876,
    10/18.1939/19.3255,
    11/18.3793/19.5109,
    12/18.5711/19.7027,
    14/18.4960/19.6276,
    15/18.1995/19.3311,
    16/18.3478/19.4794,
    17/18.4021/19.5337,
    18/18.4679/19.5995,
    19/18.6780/19.8096,
    20/18.8602/19.9917
  } {
    \draw[very thick,layerindigo] (\a,\y)--(\b,\y);
    \draw[layerindigo] (\a,\y-0.22)--(\a,\y+0.22);
    \draw[layerindigo] (\b,\y-0.22)--(\b,\y+0.22);
  }
  \draw[very thick,accent] (17.8880,13)--(19.0196,13);
  \draw[accent] (17.8880,12.7)--(17.8880,13.3);
  \draw[accent] (19.0196,12.7)--(19.0196,13.3);
  \node[accent,anchor=west,font=\scriptsize] at (17.82,0.62) {one interval misses};
  \foreach \t in {18,19,20} {\draw[layergrey] (\t,-0.2)--(\t,0.2);}
  \foreach \t in {18,19,20} {\node[below,font=\scriptsize] at (\t,-0.45) {\t};}
\end{tikzpicture}
\caption{Twenty $95$ percent confidence intervals for $\mu$, each from one
sample of size $n=12$ with $\sigma=1$ known, so the half-width is
$1.96/\sqrt{12}=0.5658$ and the width is $1.1316$. The vertical dashed line is the
true $\mu=19.125$. Nineteen intervals contain it; the interval in row
13, namely $(17.8880;\,19.0196)$, does not. One failure in twenty is a
reasonable number for a $95$ percent promise, not a sign of a broken procedure.}
\label{fig:dua-puluh-selang}
\end{figure}

\begin{catatan}
Three sentences that are true, and one that is false. True: ``this procedure catches $\mu$
on $95$ percent of repeated trials''. True: ``we are $95$ percent confident that $\mu$
lies between $18.44$ and $19.81$'', provided ``confident'' is read as a summary of
repeated trials. True: ``this procedure is wrong in about five cases out of a hundred''.
False: ``the probability that $\mu$ is in this interval is $0.95$''. The last sentence
pretends to know something about $\mu$ that the data do not provide.
\end{catatan}

The width of an interval shrinks like $1/\sqrt n$, not $1/n$: halving it demands four times
the data, making it a tenth demands a hundred times. This rule can be computed before the
experiment starts, since a half-width $z\sigma/\sqrt n\le d$ is equivalent
to $n\ge(z\sigma/d)^2$; with $\sigma=1$, $z=1.96$, and $d=0.1$, we need
$n\ge385$.

\begin{tangan}
Prepare twenty cards, write the numbers $1$ to $20$ on the cards, one number per card.
Shuffle, draw five cards without looking, add them up, then divide by five. Write the result on
one line of paper. Return the five cards, shuffle again, repeat. Do this thirty times, always
five cards, and arrange those numbers as points on one number line.

What you should see: the points do not pile up at one number. The mean of all the
cards is really $10.5$, and your points scatter around it, some $8$, some
$13$. Now repeat the whole job with ten cards per draw, thirty
times more. The second spread is narrower, but not by half: it narrows by roughly
the square root of two, because the number of cards doubled. Paste the two rows of points side
by side, and see for yourself that doubling the work only trims the wobble by about thirty
percent.
\end{tangan}

\section{Proportions: Wald, Wilson, and the rule of three}
\tingkat{intermediate}

For proportions, the most widely used formula is the simplest version of the previous
theorem: replace $\bar X_n$ by $\hat p$ and $\sigma/\sqrt n$ by the standard error of the
proportion. In everyday language: we set a symmetric interval around the observed proportion,
twice its standard error wide.

\begin{definisi}[Wald interval]
\index{wald interval}
The Wald interval of level $1-\alpha$ for a proportion is
\[
\hat p\pm z\sqrt{\frac{\hat p(1-\hat p)}{n}},\qquad z=z_{1-\alpha/2}.
\]
\end{definisi}

That formula is reasonable as long as $\hat p$ is far from zero and one. The difficulty
appears exactly where the data are scarcest, and its shape is not smooth. If $\hat p=0$, the
formula gives the interval $[0;\,0]$, of zero width: a claim of complete confidence that the
event rate is zero, when all we have is the absence of events in $n$ trials. If $\hat p=1/20$,
the lower limit is $-0.0455$, a number that is not a probability.

The cause is not rounding. Wald uses $\hat p$ for two jobs at once: as the center point and as
the variance estimator. When $\hat p$ approaches zero, the estimated variance shrinks with it,
although the real uncertainty is precisely large on that scale.

The way out is to invert the question: instead of assuming that $\hat p$ spreads normally
around the true $p$, we ask which values of $p$ are still plausible given the observed data.
The criterion is that the normalized distance between $\hat p$ and $p$ does not exceed $z$,
with the variance in the denominator using $p$, not $\hat p$.

\kamus{Wilson interval}{interval wilson}{an interval found by asking which values are still
plausible, rather than by setting a symmetric interval around the guess}

\begin{teorema}[Wilson interval]
\index{wilson interval}
For $0\le k\le n$, $\hat p=k/n$, and $z>0$, the set of $p\in[0,1]$ satisfying
\[
n(\hat p-p)^2\le z^2p(1-p)
\]
is a closed interval with endpoints
\[
\frac{\hat p+\dfrac{z^2}{2n}\ \pm\
z\sqrt{\dfrac{\hat p(1-\hat p)}{n}+\dfrac{z^2}{4n^2}}}{1+\dfrac{z^2}{n}} .
\]
Both endpoints always lie in $[0,1]$, and the interval is never empty.
\end{teorema}

\begin{proof}[Proof]
Rearrange the inequality into $(n+z^2)p^2-(2n\hat p+z^2)p+n\hat p^2\le0$.
Its discriminant is
\[
(2n\hat p+z^2)^2-4(n+z^2)n\hat p^2
=4z^2\Bigl(n\hat p(1-\hat p)+\frac{z^2}{4}\Bigr)\ge0,
\]
so the roots are real and the quadratic formula gives the form in the theorem statement after
numerator and denominator are divided by $2n$.

For the location, let $f(p)=n(\hat p-p)^2-z^2p(1-p)$. The coefficient of $p^2$ is
$n+z^2>0$, so $\{f\le0\}$ is the interval $[r_-,r_+]$. Because
$f(0)=n\hat p^2\ge0$, the point $0$ is not in the interior of that interval, so
$r_-\ge0$; because $f(1)=n(1-\hat p)^2\ge0$, neither is the point $1$, so
$r_+\le1$. Finally $f(\hat p)=-z^2\hat p(1-\hat p)\le0$ places $\hat p$ inside the
interval, so the interval is not empty.
\end{proof}

Now compare the two on the same data. The following table uses $z=1.96$ and
$95$ percent confidence intervals; every number is computed with the formulas above and
rechecked with Python.

\begin{center}
\small
\begin{tabular}{lrrrr}
\hline
data & $\hat p$ & Wald $95$ percent & Wilson $95$ percent & Wilson width \\
\hline
$33/50$  & 0,660 & $(0.5287;\ 0.7913)$ & $(0.5215;\ 0.7756)$ & 0,254 \\
$53/100$ & 0,530 & $(0.4322;\ 0.6278)$ & $(0.4329;\ 0.6249)$ & 0,192 \\
$13/20$  & 0,650 & $(0.4410;\ 0.8590)$ & $(0.4329;\ 0.8188)$ & 0,386 \\
$1/20$   & 0,050 & $(-0.0455;\ 0.1455)$ & $(0.0089;\ 0.2361)$ & 0,227 \\
$0/20$   & 0,000 & $(0;\ 0)$ & $(0;\ 0.1611)$ & 0,161 \\
\hline
\end{tabular}
\end{center}

The first two rows barely differ: when $\hat p$ is near one half, both are legitimate and
agree to two decimal places. The fourth and fifth rows show the separation. At $1/20$ the Wald
interval is \emph{narrower} than Wilson, $0.191$ against $0.227$, and that narrowness is
the sign of an error: its lower limit falls outside $[0,1]$. At $0/20$ Wald gives zero width,
whereas Wilson leaves room up to $0.1611$. The conclusion is not that Wilson is always
better, but that the width of an interval cannot be judged without its coverage.

For the zero-event case there is a special result older than Wilson and easier to remember. In
everyday language: no events at all is not proof that the event rate is zero, and the upper
limit is about three divided by the number of trials.

\kamus{rule of three}{aturan tiga}{if no event occurs in a given number of trials,
the event rate may still be as high as three divided by the number of trials}

\begin{teorema}[Rule of three]
\index{rule of three}
Let $K\sim\mathrm{Bin}(n,p)$ and suppose $K=0$ is observed. The one-sided upper limit at the
$95$ percent level is $p_{\mathrm{hi}}=1-0.05^{1/n}$, and
\[
p_{\mathrm{hi}}\le\frac{\ln 20}{n}<\frac{3}{n}.
\]
\end{teorema}

\begin{proof}[Proof]
The probability of observing zero events is $(1-p)^n$. Setting that probability equal to
$0.05$ gives $(1-p_{\mathrm{hi}})^n=0.05$, that is
$p_{\mathrm{hi}}=1-0.05^{1/n}=1-e^{-c/n}$ with $c=\ln 20=2.9957\ldots$ Because
$1-e^{-x}\le x$ for $x\ge0$, we obtain $p_{\mathrm{hi}}\le c/n<3/n$.
\end{proof}

For $n=100$, the rule of three gives $3/100=0.03$, the exact limit $0.0295130$, and the
two-sided Wilson upper limit at $k=0$ is $z^2/(n+z^2)=0.0369935$. All three numbers answer
the same question with three different promises; the rule of three is the exact one-sided
limit, Wilson a two-sided approximation. What is not allowed is to quote a number without
naming the procedure.

\section{Hypothesis tests and p-values}
\tingkat{intermediate}

Up to here the question took the form ``what is its value''. The next question takes the form
``is its value equal to a particular number''. In everyday language, we file a charge that
there is no effect, then check whether the data are surprising enough to overturn that charge.

\begin{definisi}[Null hypothesis, alternative, and level]
\index{null hypothesis}
The \emph{null hypothesis} $H_0$ is a statement about the parameter to be tested; the
\emph{alternative hypothesis} $H_1$ is the set of values held to be true if $H_0$ is
rejected. The \emph{level} $\alpha$ is the probability of rejecting $H_0$ when $H_0$ is in fact
true. The rejection rule is determined by one test statistic $T$ and one rejection region
$R$: we reject $H_0$ if $T\in R$.
\end{definisi}

The level $\alpha$ is chosen before the data are seen, and therein lies the whole credibility
of the procedure: we are willing to wrongly accuse in a small fraction of trials, and we state
how large that fraction is. The p-value moves the same idea to a single data set.

\begin{definisi}[p-value]
\index{p-value}
For data $x$ and test statistic $T$, the p-value is the probability of obtaining a result at
least as extreme as the observed one if $H_0$ is true,
\[
\text{p-value} p=\sup_{\theta\in H_0}P_\theta\bigl(T\text{ as far as }T(x)\bigr).
\]
We reject $H_0$ at level $\alpha$ if the p-value $\le\alpha$.
\end{definisi}

The phrase ``at least as far'' needs to be fixed for a two-sided test, and the usual
convention is to add the two tails: for a symmetric $z$ statistic, the
p-value $=2(1-\Phi(|z|))$. That definition makes the p-value a random
variable, and its behavior under $H_0$ can be stated exactly.

\begin{teorema}[The p-value is uniform under $H_0$]
\index{p-value}
If $H_0$ is true and $T$ has a continuous distribution with distribution function $F_0$, then
$P_{H_0}(\text{p-value} p\le\alpha)=\alpha$ for every $\alpha\in(0,1)$. Consequently
$P_{H_0}(\text{p-value} p\le0.05)=0.05$: a test with the rule ``reject if the p-value
$\le\alpha$'' has level exactly $\alpha$.
\end{teorema}

\begin{proof}[Proof]
For a right one-sided test, the p-value $=1-F_0(T)$, so
\[
P_{H_0}(1-F_0(T)\le\alpha)
=P_{H_0}\bigl(F_0(T)\ge1-\alpha\bigr)
=P_{H_0}\bigl(T\ge F_0^{-1}(1-\alpha)\bigr)=\alpha,
\]
because $F_0(T)$ is uniformly distributed on $(0,1)$ when $F_0$ is continuous. For a two-sided
test, the p-value depends on $|T|$ and the proof uses the same step on the
distribution of $|T|$.
\end{proof}

That uniformity has a consequence often missed: under $H_0$ the p-value spreads evenly, so a
p-value $=0.04$ appears in four percent of trials whose null hypothesis is true. What
distinguishes two trials is not the single p-value but the power of each.

\begin{bacakeras}
Two voices, taking turns.

``The p-value is point zero four. So the probability that there is no effect is four
percent.''

``No. The p-value counts the data, not the charge. If there were no effect, and
this trial were repeated a hundred times, about four times the data would look this far out.''

``So it is the data that get counted.''

``Yes, the data. The charge is only a place to stand.''

``And if the p-value is large?''

``Then the data are not surprising. The silence of a scale is not proof that the goods are
light.''

Say it once more, slowly: the p-value is the probability of the data, not the probability of
the charge. Repeat that line until it feels natural in the mouth: probability of the data, not
probability of the charge.
\end{bacakeras}


\begin{resep}
\begin{enumerate}
  \item Write $H_0$ and $H_1$ as statements about one parameter, before
        computing anything.
  \item Fix $\alpha$ and its rejection region, then pause. Do not look at
        the data yet.
  \item Compute the test statistic with the standard error under $H_0$, then the p-value.
  \item Compare the p-value with the $\alpha$ already fixed, and report the
        interval together with that decision.
\end{enumerate}
A warning about order: step 2 must be finished before step 3. If $\alpha$ is chosen
after the data are seen, $\alpha$ is no longer a promise but a record of the result, and step 4
loses its meaning. What may not be swapped: step 2 and step 3.
\end{resep}

Now an example with numbers. For the data $33$ out of $50$ and $H_0\colon p=0.5$,
the $z$ statistic uses the standard error under $H_0$, namely $\sqrt{0.25/50}=0.070711$:
\[
z=\frac{0.66-0.5}{0.070711}=2.2627,\qquad
\text{p-value} p=2(1-\Phi(2.2627))=0.0237 .
\]
The exact binomial test on the same rule, $P(X\ge33)+P(X\le17)$ for
$X\sim\mathrm{Bin}(50;0.5)$, gives $0.0328$. Both numbers reject $H_0$ at the five percent
level, but the normal approximation reports a p-value almost thirty percent smaller. At $n=50$
a difference that large is enough to flip the decision when the level is $0.03$; a good
report states which test was used, not only its p-value.

Figure~\ref{fig:dua-ekor} places those two limits on the distribution of $z$ under $H_0$: two
tails of area $0.0250$ each, and the observed statistic $2.2627$ inside the right tail.

\begin{figure}[htbp]
\centering
\begin{tikzpicture}
\begin{axis}[
  width=11.2cm, height=5.4cm,
  xmin=-4, xmax=4, ymin=0, ymax=0.46,
  axis lines=left,
  xtick={-1.96,0,1.96},
  xticklabels={$-1.96$,$0$,$1.96$},
  ytick=\empty,
  tick label style={font=\small},
  xlabel={value of $z$ if $H_0$ is true},
  xlabel style={font=\small},
]
\addplot[fill=accent!25,draw=none,domain=-3.95:-1.96,samples=100]
  {0.39894228*exp(-0.5*x^2)} \closedcycle;
\addplot[fill=accent!25,draw=none,domain=1.96:3.95,samples=100]
  {0.39894228*exp(-0.5*x^2)} \closedcycle;
\addplot[very thick,layerindigo,domain=-3.95:3.95,samples=200]
  {0.39894228*exp(-0.5*x^2)};
\draw[densely dashed,layergrey] (axis cs:-1.96,0)--(axis cs:-1.96,0.0584);
\draw[densely dashed,layergrey] (axis cs:1.96,0)--(axis cs:1.96,0.0584);
\draw[very thick,accent] (axis cs:2.2627,0)--(axis cs:2.2627,0.0308);
\node[font=\scriptsize] (tl) at (axis cs:-2.55,0.135) {$0.025$};
\draw[-{Stealth[length=4pt]},layergrey] (tl)--(axis cs:-2.55,0.012);
\node[accent,font=\scriptsize] (tz) at (axis cs:2.05,0.215) {$z=2.2627$};
\draw[-{Stealth[length=4pt]},accent] (tz)--(axis cs:2.2627,0.055);
\node[font=\scriptsize] (tr) at (axis cs:3.05,0.135) {$0.025$};
\draw[-{Stealth[length=4pt]},layergrey] (tr)--(axis cs:3.05,0.004);
\end{axis}
\end{tikzpicture}
\caption{The distribution of the statistic $z=\sqrt n(\hat p-p_0)/\sqrt{p_0(1-p_0)}$ when
$H_0$ is true, for the data $33$ out of $50$ with $H_0\colon p=0.5$. The two dashed lines at
$\pm1.96$ cut off the two shaded tails, each of area
$1-\Phi(1.96)=0.0250$, so the two-sided level is $0.05$. The thick line at
$2.2627$ is the statistic observed on that data; the right tail from that point has area
$0.0118$, and twice that area is the p-value $0.0237$ that the normal test reports.}
\label{fig:dua-ekor}
\end{figure}

A confidence interval and a hypothesis test are not two different procedures; they are two
ways of reading the same set.

\begin{proposisi}[Equivalence of interval and test]
\index{confidence interval}
For the Wald interval and the two-sided $z$-test that uses the standard error $\hat p$, we have
that $p_0$ is contained in the Wald interval of level $1-\alpha$ if and only if $H_0\colon
p=p_0$ is not rejected at level $\alpha$. More generally, the Wilson interval is the set of
$p_0$ not rejected by the score test at level $\alpha$.
\end{proposisi}

\begin{proof}[Proof]
For the Wald pair, both statements use the same inequality, namely
$|\hat p-p_0|\le z\sqrt{\hat p(1-\hat p)/n}$: the first reads it as
``$p_0$ lies inside the interval'', the second as ``the test statistic does not exceed
$z$''. For Wilson, the set $\{p_0:n(\hat p-p_0)^2\le z^2p_0(1-p_0)\}$ is the
definition of the score test turned around, and the Wilson theorem has already computed its
endpoints.
\end{proof}

\section{Type I error, type II error, and power}
\tingkat{advanced}

Every binary decision can be wrong in two ways, and the two are not symmetric.
In everyday language: a fire alarm can sound when there is no fire, and
can stay silent when there is one. The first mistake sends us out of the house for nothing;
the second is far more costly.

\begin{definisi}[Type I error, type II error, and power]
\index{type i error}
A \emph{type I error} is rejecting $H_0$ when $H_0$ is true; its probability is the level
$\alpha$. A \emph{type II error} is not rejecting $H_0$ when $H_1$ is true;
its probability is written $\beta$. The \emph{power} of a test is $1-\beta$, the probability of
rejecting $H_0$ when the effect sought is really there.
\end{definisi}

Figure~\ref{fig:dua-galat} arranges those four possibilities as a decision box,
with probabilities from the exact computation at $n=20$, $p_0=0.5$, and $p_1=0.75$ that
appears in the computational lab: $\alpha=0.0414$, power $0.6172$, so $\beta=0.3828$.

\begin{figure}[htbp]
\centering
\begin{tikzpicture}[font=\small,
  keputusan/.style={draw=layergrey,thick,rounded corners=2pt,
    minimum width=4.3cm,minimum height=1.6cm,align=center,inner sep=4pt}]
\node[keputusan,fill=layerjade!10] (a) at (0,0)
  {does not reject $H_0$\\ correct, probability $1-\alpha=0.9586$};
\node[keputusan,fill=accent!12] (b) at (5.1,0)
  {rejects $H_0$\\ type I error, probability $\alpha=0.0414$};
\node[keputusan,fill=accent!12] (c) at (0,-2.4)
  {does not reject $H_0$\\ type II error, probability $\beta=0.3828$};
\node[keputusan,fill=layerjade!10] (d) at (5.1,-2.4)
  {rejects $H_0$\\ correct, power $1-\beta=0.6172$};
\node[font=\scriptsize,above=2pt] at (0,1.0) {$H_0$ is true: $p=0.5$};
\node[font=\scriptsize,above=2pt] at (5.1,1.0) {$H_1$ is true: $p=0.75$};
\end{tikzpicture}
\caption{The two kinds of error as four decision boxes. Top row: the actual
state $H_0$ is true; bottom row: $H_1$ is true with $p=0.75$. Left column: the test
does not reject; right column: the test rejects. The boxes in red are the two errors that
cannot be suppressed together: shrinking $\alpha=0.0414$ enlarges
$\beta=0.3828$, and conversely. The numbers are computed from the exact binomial
distribution at $n=20$ with the rejection rule $k\le5$ or $k\ge15$.}
\label{fig:dua-galat}
\end{figure}

Power is not a property of the test alone, but of the test together with the sample size and
the size of the effect. Hence it must be computed for one value $p_1$ deemed important, not in
general. For a two-sided test with $H_0\colon p=0.5$, $n=100$, and the rule to reject when
$|z|>1.96$, the rejection region is $k\le40$ or $k\ge60$. The exact binomial computation
gives size $0.0569$, not $0.05$, and power $0.8750$ at
$p_1=0.65$, whereas the normal approximation gives $0.8508$. The normal approximation,
namely
\[
1-\beta\approx\Phi\Bigl(\frac{(p_1-p_0)\sqrt n}{\sqrt{p_0(1-p_0)}}-z\Bigr),
\]
is close enough for planning, but the two-and-a-half-point difference is real and comes from
the discreteness of the data.

First lesson: the nominal level is seldom attained exactly. At $n=20$ the same rule rejects
when $k\le5$ or $k\ge15$, with exact size $0.0414$; at $n=100$
the size is $0.0569$. A discrete test cannot put a threshold between two whole
numbers, so the actual level hangs on $n$ and $p_0$. Second lesson: at
$n=20$ the power against $p_1=0.75$ is only $0.6172$, so almost four in ten
trials of size twenty fail to find it.

That last number opens the most costly misinterpretation: reading ``p-value
$>0.05$'' as ``no effect''. First, the p-value is
$P(\text{data}\mid H_0)$, not $P(H_0\mid\text{data})$; both are conditional probabilities
with the conditions swapped, and we already met that swap in the diagnostic tests of
Chapter 17. Second, a large p-value is a statement about the data, not about the world: scanty
data cannot surprise anyone. Third, failure to reject can occur precisely because the test is
weak, as at $n=20$ above.

A small example makes all three concrete. Ten successes out of twenty trials against
$H_0\colon p=0.5$ give a two-sided p-value of exactly $1$: that data set sits
precisely in the middle. Is that proof that $p=0.5$? No. The $95$ percent Wilson interval
for the same data is $(0.2993;\ 0.7007)$. So the true rate may be
$0.30$, may be $0.70$, and this data set cannot tell it apart from $0.50$.
A p-value $=1$ only means ``no direction of deviation'', and that is far less than what people
usually read into it.

\begin{tantangan}
Two versions of a detector are tested on \mbox{$100\,000$} examples each. Version A catches
\mbox{$66\,000$}, version B catches \mbox{$65\,000$}. A two-proportion test gives a p-value of about
$0.0000026$, and the team writes in the report: ``A is better, $p<0.001$''. Does that
sentence have the support of the data? Answer with one interval, not with a p-value.

Hint 1: the p-value measures surprise, not the size of the effect. Compute the $95$ percent
interval for the difference of the two proportions, and mind the units.

Hint 2: the standard error of the difference is $0.002126$, so the interval runs from
$0.0058$ to $0.0142$, that is zero point six to one point four percentage points.
Compare it with the practical threshold the team itself uses: three percentage points.

If you get stuck, go back to the section Type I error, type II error, and power.
\end{tantangan}

\begin{untukanak}
A line judge is tasked with raising the flag when the ball goes out of the field. He raises the
flag only when he sees the ball go out himself, clearly enough. If he does not raise the flag,
does that mean the ball was always in? Certainly not. The ball may have gone out, but his view
was blocked by another player, or the ball was very fast.

The way he works is similar to the way a hypothesis test works. He is given one rule that must
be settled before the match: how sure he must be before raising the flag. If the rule is too
loose, he often raises the flag when the ball is in, and the crowd is angry. If the rule is too
strict, he almost never raises the flag, and big fouls slip by. The two mistakes are
different, and he cannot suppress both at once by carefulness alone.

In this chapter we shall learn why ``no evidence'' is not the same as ``no foul'', why the
number of observations matters as much as the rule chosen, and why a judge who never raises
the flag is not a more careful judge.
\end{untukanak}

\section{Worked example: 33 out of 50 and 53 out of 100}
\tingkat{intermediate}

We return to the two detectors in the chapter opening and compute both intervals all the way.
Every number below uses $z=1.96$, so $z^2=3.8416$.

For A, $\hat p=0.66$ and $n=50$. The first term inside the Wilson root is
$0.66\cdot0.34/50=0.2244/50=0.004488$; the second term
$z^2/(4n^2)=3.8416/10000=0.00038416$; their sum is $0.00487216$ and its root
$0.069801$. The center is
\[
\frac{0.66+3.8416/100}{1+3.8416/50}
=\frac{0.698416}{1.076832}=0.648584,
\]
and the half-width is $1.96\cdot0.069801/1.076832=0.127048$. So the Wilson interval
for A is $(0.5215;\ 0.7756)$, of width $0.2541$. The Wald interval for the
same data uses standard error $0.066993$ and half-width
$1.96\cdot0.066993=0.131303$, hence $(0.5287;\ 0.7913)$ of width
$0.2626$.

For B, $\hat p=0.53$ and $n=100$. The first term
$0.53\cdot0.47/100=0.002491$, the second term $3.8416/40000=0.00009604$,
the sum $0.00258704$ and its root $0.050863$. The center
$(0.53+0.019208)/1.038416=0.528890$, the half-width
$1.96\cdot0.050863/1.038416=0.096003$, hence $(0.4329;\ 0.6249)$ of width
$0.1920$. Wald gives $(0.4322;\ 0.6278)$ of width $0.1956$.

Now compare. The intervals A and B overlap on $(0.5215;\ 0.6249)$, an extent of
$0.1034$: about $41$ percent of the length of A and $54$ percent of the length of B. The
normal-based two-proportion test gives $z=1.5175$ and a two-sided p-value of $0.1291$.
Although the raw numbers differ by thirteen points, this data set is not enough to conclude
that A is better; the interval for A is also wider, because it is built from fifty examples.

\kertas{Recompute part A without looking at the text: $0.66\cdot0.34$, divide by $50$,
add $3.8416/10000$, take its root, then divide by $1.076832$ after
multiplying by $1.96$. Check that the result is $0.127048$. Then answer: how many
examples are needed for the width of the interval for A to fall to $0.15$, if $\hat p$
stays $0.66$? Hint: the width is proportional to $1/\sqrt n$, and you already know the width
at $n=50$.}

\section{Computational lab: coverage, type I error, and power}
\tingkat{advanced}

Every claim above can be checked with a computer, and there is one claim that can be checked
only that way: the actual coverage of a procedure. The key is knowing the true $p$, something
that never happens with real data. Hence we work with simulation and with exact binomial
computations.

\begin{lab}
\begin{verbatim}
import numpy as np
from math import sqrt
from scipy.stats import norm, binom

z = norm.ppf(0.975)


def wald(k, n):
    p = k / n
    se = sqrt(p * (1 - p) / n)
    return p - z * se, p + z * se


def wilson(k, n):
    p = k / n
    z2 = z * z
    mid = (p + z2 / (2 * n)) / (1 + z2 / n)
    half = z * sqrt(p * (1 - p) / n + z2 / (4 * n * n)) / (1 + z2 / n)
    return mid - half, mid + half


# 1. width of the 95% interval for several proportions
print("== 1. width of the 95% interval ==")
print("   n    k   phat  Wald width    Wilson width")
for n in (20, 100):
    for p in (0.00, 0.05, 0.25, 0.50, 0.66):
        k = round(p * n)
        a = wald(k, n)
        b = wilson(k, n)
        print("%4d %4d  %5.2f   %9.4f    %9.4f"
              % (n, k, k / n, a[1] - a[0], b[1] - b[0]))

# 2. exact coverage at p = 0.05
print("== 2. exact coverage, p=0.05 ==")
for n in (20, 50, 100):
    for nama, f in (("Wald", wald), ("Wilson", wilson)):
        c = 0.0
        for k in range(n + 1):
            lo, hi = f(k, n)
            if lo <= 0.05 <= hi:
                c += binom.pmf(k, n, 0.05)
        print("n=%3d %-6s coverage=%.4f" % (n, nama, c))

# 3. simulated coverage with a fixed seed
rng = np.random.default_rng(20261024)
R = 20000
print("== 3. simulated coverage, R=%d ==" % R)
for n in (20, 100):
    X = rng.binomial(n, 0.05, R)
    for nama, f in (("Wald", wald), ("Wilson", wilson)):
        ok = 0
        w = 0.0
        for k in X:
            lo, hi = f(int(k), n)
            ok += (lo <= 0.05 <= hi)
            w += hi - lo
        print("n=%3d %-6s coverage=%.4f  width=%.4f" % (n, nama, ok / R, w / R))

# 4. type I error and power for H0: p=0.5
print("== 4. type I error and power, H0 p=0.5 ==")


def tolak(k, n, p0=0.5):
    return abs((k / n - p0) / sqrt(p0 * (1 - p0) / n)) > z


for n in (20, 100):
    rej = [k for k in range(n + 1) if tolak(k, n)]
    lo = max(k for k in rej if k < n / 2)
    hi = min(k for k in rej if k > n / 2)
    print("n=%3d reject if k<=%d or k>=%d" % (n, lo, hi))
    print("      exact size at p=0.5 : %.4f"
          % sum(binom.pmf(k, n, 0.5) for k in rej))
    for p1 in (0.65, 0.75):
        ex = sum(binom.pmf(k, n, p1) for k in rej)
        zc = (p1 - 0.5) / sqrt(0.25 / n)
        ap = norm.cdf(zc - z) + norm.cdf(-zc - z)
        print("      power at p=%.2f: exact %.4f  approx %.4f"
              % (p1, ex, ap))

# 4b. simulated type I error rate when H0 is true
print("== 4b. simulated type I error, R=%d ==" % R)
for n in (20, 100):
    X0 = rng.binomial(n, 0.5, R)
    tolak0 = np.abs((X0 / n - 0.5) / sqrt(0.25 / n)) > z
    print("n=%3d p=0.5: reject=%.4f" % (n, tolak0.mean()))

# 5. t-interval for 12 measurements
x = np.array([18.2, 19.1, 17.8, 20.3, 18.9, 19.7,
              18.4, 21.0, 17.5, 19.9, 18.6, 20.1])
n = len(x)
m = x.mean()
s = x.std(ddof=1)
from scipy.stats import t
tq = t.ppf(0.975, n - 1)
h = tq * s / sqrt(n)
print("== 5. t-interval, n=12 ==")
print("mean=%.6f  s=%.6f  se=%.6f" % (m, s, s / sqrt(n)))
print("t=%.6f  half-width=%.6f  CI=(%.6f, %.6f)" % (tq, h, m - h, m + h))
print("half-width with z=1.96: %.6f" % (z * s / sqrt(n)))
\end{verbatim}
\end{lab}

The actual output, verbatim:

\begin{verbatim}
== 1. width of the 95% interval ==
   n    k   phat  Wald width    Wilson width
  20    0   0.00      0.0000       0.1611
  20    1   0.05      0.1910       0.2272
  20    5   0.25      0.3795       0.3568
  20   10   0.50      0.4383       0.4014
  20   13   0.65      0.4181       0.3860
 100    0   0.00      0.0000       0.0370
 100    5   0.05      0.0854       0.0902
 100   25   0.25      0.1697       0.1676
 100   50   0.50      0.1960       0.1923
 100   66   0.66      0.1857       0.1826
== 2. exact coverage, p=0.05 ==
n= 20 Wald   coverage=0.6389
n= 20 Wilson coverage=0.9245
n= 50 Wald   coverage=0.9199
n= 50 Wilson coverage=0.9622
n=100 Wald   coverage=0.8775
n=100 Wilson coverage=0.9659
== 3. simulated coverage, R=20000 ==
n= 20 Wald   coverage=0.6430  width=0.1472
n= 20 Wilson coverage=0.9224  width=0.2193
n=100 Wald   coverage=0.8783  width=0.0829
n=100 Wilson coverage=0.9667  width=0.0884
== 4. type I error and power, H0 p=0.5 ==
n= 20 reject if k<=5 or k>=15
      exact size at p=0.5 : 0.0414
      power at p=0.65: exact 0.2457  approx 0.2687
      power at p=0.75: exact 0.6172  approx 0.6088
n=100 reject if k<=40 or k>=60
      exact size at p=0.5 : 0.0569
      power at p=0.65: exact 0.8750  approx 0.8508
      power at p=0.75: exact 0.9997  approx 0.9988
== 4b. simulated type I error, R=20000 ==
n= 20 p=0.5: reject=0.0430
n=100 p=0.5: reject=0.0554
== 5. t-interval, n=12 ==
mean=19.125000  s=1.082191  se=0.312401
t=2.200985  half-width=0.687591  CI=(18.437409, 19.812591)
half-width with z=1.96: 0.612296
\end{verbatim}

Four things deserve reading from that output. First, parts 2 and 3 agree to almost
three digits, and that agreement is itself a check, since a simulation with
$R=20000$ repetitions can approach the exact computation only to about
$\sqrt{0.25/20000}=0.0035$. At $p=0.05$ and $n=20$ the difference is $0.0041$ for
Wald and $0.0021$ for Wilson. An honest simulation must report $R$.

Second, the coverage of Wald at $n=20$ and $p=0.05$ is $0.6389$, not $0.95$.
The cause can be computed by hand: the probability of observing zero events is $0.3585$,
and at $k=0$ the Wald interval is $[0;0]$, which does not contain $0.05$; that single point
alone consumes almost thirty-six points of coverage. Wilson fails only when
$k\ge3$, whose probability is $0.0755$, and its failure is not total because its interval
stays wide. That is why it reaches $0.9245$. Note also the width column: the mean
Wald width $0.1472$ is smaller than the Wilson $0.2193$, and precisely therein lies the
point. A narrower interval is not a better interval if it is narrow because it sits in the
wrong place.

Third, part 4 separates the size of a test from the nominal level: $0.0414$ at $n=20$ and
$0.0569$ at $n=100$, both for the rule $|z|>1.96$; part 4b repeats it
by simulation and lands close, $0.0430$ and $0.0554$. Power moves far more
drastically: $0.2457$ at $n=20$ and $0.8750$ at $n=100$ for $p_1=0.65$.
The normal approximation is not always on the safe side: at $n=20$ it reports $0.2687$
while the truth is $0.2457$.

Fourth, part 5 merely repeats the $t$-interval already computed in the text, with a
difference in half-width between $t$ and $z$ of $0.075295$.

The program does not prove any one of all those things. It checks claims already
proved, and in two places it shows its own limits: exact coverage
uses the distribution we assume to be true, so if the model is wrong, the number
$0.9245$ means nothing; and simulation only measures how often
the procedure succeeds under that model, not how right the model is.

\begin{untukdosen}
The teaching order I suggest: start from the morning meeting in the chapter opening, not from
the definition of an estimator. Write $33/50$ and $53/100$ on the board, ask the class to choose, then
ask what happens if those fifty examples are replaced. Out of that unease the
definitions of estimator, bias, and standard error emerge as the answer. The MSE
decomposition is proved at the board with one line of algebra. The interpretation of the
confidence interval comes after that, as material for a lecture of its own, not a footnote.

Two misconceptions appear most often. First, students read ``$95$ percent''
as a conditional probability on data already in hand, then write ``the probability that $\mu$ is in
this interval is $0.95$''; the symptom shows when they cannot explain what is
random about an interval already computed. Second, they take a large p-value as
proof that $H_0$ is true, and will compute the p-value for $10/20$ and then conclude that
$p=0.5$ is established.

A blackboard activity that works: split the class into two groups, one computing the
Wald interval and one the Wilson interval for $1/20$, then have both defend their result
in front of the class. The Wald group will find a negative lower limit and have to
explain it. The sign that the class is ready to go on: they can name one example of a
sample size that makes the p-value large although the effect is large, and can explain
why that happens.
\end{untukdosen}

\section{Summary}
\tingkat{basic}

\begin{ringkasan}
\begin{itemize}
  \item An estimator is a statistic that guesses a parameter, so it is a random variable with
        an expectation and a variance. Its quality is measured by
        $\operatorname{MSE}=\operatorname{bias}^2+\operatorname{Var}$; unbiasedness alone
        guarantees nothing, and MSE going to zero guarantees consistency through
        Chebyshev.
  \item The standard error is the square root of the variance of the estimator, not the width
        of the data. For a proportion its value is $\sqrt{\hat p(1-\hat p)/n}$, shrinking like
        $1/\sqrt n$, so four times the data only cuts it in half.
  \item A $95$ percent confidence interval is a statement about the procedure, not
        about the parameter: the procedure catches the true value on $95$ percent
        of repeated trials. Once the data are taken, the interval is fixed and no longer
        carries a probability.
  \item The Wald interval collapses near $0$ and $1$: at $0/20$ its width is zero, and at
        $p=0.05$ with $n=20$ its coverage is only $0.6389$. The Wilson interval, the
        result of inverting the score test, reaches $0.9245$ in the same setting.
  \item If no event occurs in $n$ trials, the $95$ percent upper limit is about
        $3/n$; for $n=100$ the exact limit is $0.0295130$ and the two-sided Wilson
        limit $0.0369935$.
  \item The p-value is the probability of data at least that extreme if $H_0$ is true, and
        under $H_0$ it is uniformly distributed. A large p-value is not proof that $H_0$ is
        true: at $n=20$ the power against $p=0.75$ is only $0.6172$.
\end{itemize}
\end{ringkasan}

\section{Exercises}
\tingkat{intermediate}

\begin{latihan}
\begin{enumerate}
  \item Assess this claim: ``Because the p-value equals $0.42$, the null hypothesis
        $p=0.5$ is proved true.'' True, false, or not enough evidence? Name
        the part that is right, the part that is wrong, and one additional computation that
        would strengthen your answer.
  \item Prove that $\mathbb{E}[\hat p(1-\hat p)]=p(1-p)(n-1)/n$ for
        $\hat p=K/n$ with $K\sim\mathrm{Bin}(n,p)$, then determine the unbiased estimator
        of $p(1-p)$ obtained from it. Check the result at $p=0.3$ and
        $n=20$.
  \item Prove that if $\operatorname{bias}(\hat\theta_n)\to0$ and
        $\operatorname{Var}(\hat\theta_n)\to0$, then $\hat\theta_n$ is consistent for
        $\theta$. Explain why that condition is weaker than
        $\operatorname{MSE}\to0$ yet sufficient for the same conclusion.
  \item Compute the $95$ percent Wilson and Wald intervals for $7/40$, then compare
        their widths with $13/20$. Explain why the smaller width at $7/40$
        is not a sign of higher precision, and check both intervals with
        Python.
  \item For $n=30$ and $H_0\colon p=0.5$ with a normal-based two-sided test,
        determine the rejection region, compute its exact size, then compute its power at
        $p_1=0.7$. Compare that power with $n=100$ at the same
        $p_1$.
  \item Prove that the Wald interval and the two-sided $z$-test using the standard error
        $\hat p$ really are equivalent, by writing the two inequalities side by side.
        Explain why the equivalence fails when the test uses the standard error under
        $H_0$.
  \item A report writes: ``there is no significant difference, $p=0.31$,
        so the two methods are equally good''. Rewrite that sentence honestly, naming the
        interval, the level, and the power. You may invent the numbers as long as they are
        consistent.
  \item For $n=100$ with not a single event, compute the rule of three, the exact one-sided
        limit, and the Wilson upper limit. Explain why the three numbers differ, then
        state which one may be used for the sentence ``at most $5$ percent''.
  \item Compute the $90$ percent $t$-interval for the twelve delay-time data in this chapter,
        then compare it with the $95\%$ interval. How much data is needed for the width
        of the $95\%$ interval to fall by half, if $s$ does not change?
\end{enumerate}

Hints for the three hardest items. Item 3: write
$\hat\theta_n-\theta=\operatorname{bias}_n+(\hat\theta_n-\mathbb{E}\hat\theta_n)$,
then use the triangle inequality and Chebyshev on the second term. Item 6: the first inequality
is $|\hat p-p_0|\le z\sqrt{\hat p(1-\hat p)/n}$, and a test that uses
$\sqrt{p_0(1-p_0)/n}$ uses a different denominator, so the set of $p_0$ not
rejected is no longer the Wald interval but a more complicated set. Item 5: the rejection
region is $|k-15|>1.96\sqrt{30\cdot0.25}$, so $k\le9$ or $k\ge21$, and
the power is computed by summing the binomial pmf at $p=0.7$ over those two tails.
\end{latihan}

\begin{untukpeneliti}
What is really proved in this chapter: the MSE decomposition; consistency from MSE going to
zero; exact coverage of the $z$-interval under the normal model; the closed form of the Wilson
interval and its location in $[0,1]$; the rule of three as an exact one-sided limit together
with the inequality $p_{\mathrm{hi}}<\ln 20/n$; the uniformity of the p-value under $H_0$ for a
continuous statistic; and the equivalence of interval and test in the pairs mentioned.
What is merely popular, not a theorem: that $95\%$ means a conditional probability on
data already observed; that a large p-value supports $H_0$; and that a smaller interval width
is always better, when a comparison of widths without coverage means nothing. The subnominal
coverage of Wilson at some $(n,p)$ is also a computational fact, not a theorem.

Primary source: Wasserman, \emph{All of Statistics}, Chapter 6 (Models, Statistical
Inference, and Learning) for bias, MSE, and the interpretation of intervals; Chapter 9
(Parametric Inference) for intervals and the plug-in method; Chapter 10 (Hypothesis Testing
and p-values) for p-values, type I and II errors, power, and the relation of tests to
intervals. The Wilson interval comes from E.~B. Wilson, ``Probable Inference, the Law of
Succession, and Statistical Inference'' (1927); the rule of three from B.~D. Jovanovic and
P.~S. Levy, ``A Look at the Rule of Three'' (1997). Neither is discussed by Wasserman, so the
references are cited directly, not through an intermediate textbook.

Limits of validity: the whole chapter works with independent and identically distributed
samples, one parameter, and a known normal or binomial approximation. The $z$-interval is
exact only under the normal model with $\sigma$ known; beyond that it is an asymptotic
statement, and at small $n$ with heavy-tailed distributions its coverage can fall far below
nominal. The Wilson interval is a score approximation, not an exact interval; exact intervals
for proportions are computed from the binomial distribution and are wider. Every reported
coverage is conditional on the binomial model being correct.

Open questions worth a reader's work: how well proportion intervals behave when $n$ is not
fixed but stops adaptively, as in sequential sampling; and how to build intervals for
classification metrics computed from dependent data, for instance when one example appears in
several cross-validation folds. Both are still research material.
\end{untukpeneliti}

\peta{Chapter 20 in Wasserman, \emph{All of Statistics}: Chapter 6 (Models, Statistical
Inference, and Learning: estimators, bias, MSE, standard error, interpretation of intervals),
Chapter 9 (Parametric Inference: confidence intervals and the plug-in method), Chapter 10
(Hypothesis Testing and p-values: p-values, type I and II errors, power, equivalence of tests
and intervals). The Wilson interval and the rule of three are taken from their original
papers, Wilson (1927) and Jovanovic--Levy (1997), because Wasserman discusses neither. The law
of large numbers, the central limit theorem, and Chebyshev's inequality used here are proved
in Chapter 18; conditional probability and the error of swapping the conditions are discussed
in Chapter 16. Chapter 25 continues this chapter for classification metrics, including the
rule of three in the reporting of model evaluation.}
